\section{第1問：Folditと人間参加型の科学探索}
\subsection{科学的位置づけ}
Folditでは、アミノ酸鎖の立体構造を画面上で操作し、Rosettaに基づく評価で高いスコアを目指す。
人間の空間的な直観と、構造の局所最適化を組み合わせることが特徴である\cite{foldit,cooper}。
2024年ノーベル化学賞では、Bakerの計算によるタンパク質設計と、Hassabis・Jumperの構造予測が受賞した\cite{nobel}。
Folditそのものが受賞したという意味ではなく、人間参加型探索と機械による予測が同じタンパク質科学を支えるという関連として捉える。

\subsection{観測計画と実装}
\pending{著者が後日プレーする。現在の版には本人のプレー時間・スコア・順位・スクリーンショットが未収録であり、第1問は未完成である。以下は実験設計である。}
ログ収集CLIを実装し、パズルID、操作種別、操作前後のスコア、時間、順位、参加者数、画面のファイル名、自由記述を記録できるようにした。
未記録の状態で分析すると\texttt{awaiting\_real\_play}を返す。
同じ保存構造から手操作、Shake、Wiggle、Shake後のWiggleを比較する。
比較中は同じ時間を割り当て、操作順を無作為化して、習熟や疲労の影響を小さくする。
主な指標は\(\Delta S=S_{\rm after}-S_{\rm before}\)と、\(60\Delta S/t\)（1分当たりの改善）である。
異なるパズルは評価の尺度が異なるため、パズルごとに集計する。

\subsection{六つの論点を検証する方法}
直観の有効性は「疎水性部分を内側へ動かす」などの操作仮説と結果を対にして評価する。
成功例だけでなく、衝突が増えてスコアが低下した例も残す。
科学への貢献意識は、チュートリアルと研究用パズルでの意欲や操作時間の違いを本人が記述する。
共同作業・競争は、助言を受ける前後の改善と、順位を見た後に費やした時間を記録して考える。
一人の観察では競争が改善の原因とまでは断定できない。

Shakeは側鎖配置、Wiggleは連続的な構造最適化に関係するツールとして理解する\cite{foldit}。
スコアの増加だけでは局所解脱出の証拠にならない。
同じ状態からWiggleだけでは改善が止まった後、構造を変える操作を挟んでさらに高い安定スコアに達したかを確認する。
ランキングは、パズル内順位と全体Soloist順位、soloとevolverを区別する。
チュートリアルを含む総時間、最高値、最終値、計測時点の参加者数も併記する。

将来のAI応用として、人間が選んだ操作の系列を探索方策の学習データにすることが考えられる。
ただし高得点を得る操作列が、そのまま実験室で安定するタンパク質を意味するわけではない。
評価関数への過適合を避けるため、未見パズルと実験的評価の両方が必要になる。
\repolink{01}
